% ============================================================
% پایه‌های الگوریتم‌های کوانتومی و شبیه‌سازی
% سند مطالعاتی زنده: قطب‌نما، نه کتاب درسی
% ============================================================

\chapter*{این سند چگونه خوانده شود}
\addcontentsline{toc}{chapter}{این سند چگونه خوانده شود}

مخاطب این متن پژوهشگر است. هدف آن نگه‌داشتن شمای کلی است تا در پیچ‌وخم مدار و فرمول گم نشوید. از دیگر سو برای روش ها و رهیافت های اصلی توضیحات و توصیفات مفصل آورده خواهد شد تا بخش مربوطه مرجعی باشد برای فهم ریاضیات و فیزیک مسئله. هر فصل کوتاه است، شهودی است، و قابل گسترش است.

سه لایه را از هم جدا نگه دارید:
\begin{enumerate}
	\item \textbf{هدف / مسئله} - چه چیزی را می‌خواهیم بدانیم؟
	\item \textbf{رهیافت محاسباتی} - با چه مدل محاسباتی به آن می‌رسیم؟
	\item \textbf{ابزار پیاده‌سازی} - کدام زیرروال آن رهیافت را محقق می‌کند؟
\end{enumerate}
نقشه زنده در فصل~\ref{chap:map} است. اگر وسط کار پژوهشی سردرگم شدید، به همان فصل و به قطب‌نمای فصل~\ref{chap:compass} برگردید.

منبع واحد واژگان جدول~\ref{tab:glossary} در فصل~\ref{chap:glossary} است. در متن روان، فارسی مصوب می‌آید و شکل انگلیسی در نخستین کاربرد هر فصل داخل پرانتز است. سرواژه‌ها و نام الگوریتم‌ها (\lr{Trotter}، \lr{LCU}، \lr{QSP}، \lr{QPE}، \lr{VQE}) انگلیسی می‌مانند.

نسخه اول عمداً اثبات پیچیدگی، جزئیات مدار، و کد شبیه‌ساز ندارد؛ این‌ها در «جای توسعه» هر فصل فهرست شده‌اند.

% ============================================================
\part{نقشه مفهومی}
% ============================================================

\chapter{نقشه زنده}
\label{chap:map}

\section{جایگاه در نقشه}

این فصل نقشه است. بقیه فصل‌ها به اینجا ارجاع می‌دهند.

\section{ایده شهودی}

پژوهش شبیه‌سازی کوانتومی\footnote{\lr{quantum simulation}.} یک خط مستقیم از «فیزیک» به «مدار» نیست. نخست یک مسئله فیزیکی به مدل ریاضی و سپس به هامیلتونی $H$ تبدیل می‌شود. بعد باید دو چیز را جدا کرد:
\begin{itemize}
	\item \textbf{چه پرسشی از $H$ دارید؟} تحول زمانی، یا طیف / حالت پایه؟
	\item \textbf{با کدام رهیافت به آن می‌رسید؟} پیاده‌سازی $U(t)$ یا $f(H)$، مسیر آدیاباتیک، مهندسی آنالوگ، یا حلقه وردشی.
\end{itemize}
ابزارهایی مثل \lr{Trotter}، \lr{LCU} و \lr{QSP} زیرِ رهیافت‌اند، نه هم‌سطح آن. برآورد فاز\footnote{\lr{phase estimation} (\lr{QPE}).} روی مرز دینامیک و طیف می‌نشیند، چون به تحول کنترل‌شده نیاز دارد.

اگر وسط جزئیات گم شدید، به این فصل برگردید و بپرسید: سؤال من کدام است؟ رهیافت کدام است؟ ابزار کدام است؟

\section{حداقل فرمالیسم}

جریان کاری\footnote{\lr{workflow}.} کلی در شکل~\ref{fig:research-map} آمده است.

\begin{figure}[htbp]
	\centering
	\begin{latin}
			\setstretch{1}
			\resizebox{\ifdim\width>\textwidth\textwidth\else\width\fi}{!}{%
					% نقشهٔ پژوهشی: لایه‌های مسئله / پرسش / رهیافت / ابزار / اجرا
% اتصالات مطابق فصل نقشه: QPE از تحول تغذیه می‌شود و پرسش طیفی را جواب می‌دهد؛
% آنالوگ از ابزار گیتی عبور نمی‌کند؛ آدیاباتیک رهیافت است نه ورودی QPE.
% سبک مشترک نمودارهای سند پایه‌ها (فقط داخل محیط latin فراخوانی شود)
\tikzset{
	qafbox/.style={
		draw=black!55,
		rounded corners=2pt,
		align=center,
		inner sep=3pt,
		thick,
		fill=white,
		font=\footnotesize\linespread{1}\selectfont,
		minimum height=0.95cm
	},
	qafwide/.style={
		qafbox,
		text width=2.55cm
	},
	qafnarr/.style={
		qafbox,
		text width=2.05cm
	},
	qaftitle/.style={
		font=\small\bfseries\linespread{1}\selectfont,
		align=center
	},
	qafarr/.style={
		-{Latex[length=2.2mm]},
		thick,
		draw=black!65
	},
	qafpanel/.style={
		draw,
		thick,
		rounded corners=4pt,
		inner xsep=8pt,
		inner ysep=8pt
	}
}

\begin{tikzpicture}[
	font=\scriptsize\linespread{1}\selectfont,
	x=1cm,
	y=1cm
	]
	\tikzset{
		lab/.style={
			font=\tiny\bfseries\linespread{1}\selectfont,
			anchor=east,
			inner sep=1pt,
			text=black!55,
			align=right
		},
		note/.style={
			font=\tiny\linespread{1}\selectfont,
			inner sep=1pt,
			text=black!60
		}
	}

	% ----- لایه مسئله -----
	\node[qafwide, fill=blue!12] (phys) at (0,0) {Physical problem};
	\node[qafwide, fill=blue!10] (model) at (0,-1.25) {Mathematical model};
	\node[qafwide, fill=blue!14] (H) at (0,-2.5) {Hamiltonian $H$};

	% ----- لایه پرسش -----
	\node[qafnarr, fill=green!14] (dyn) at (-3.05,-4.05) {Dynamic\\question};
	\node[qafnarr, fill=orange!14] (stat) at (3.05,-4.05) {Static / spectral\\question};

	% ----- لایه رهیافت -----
	\node[qafnarr, fill=cyan!14] (ana) at (-5.15,-5.7) {Analog\\engineering};
	\node[qafnarr, fill=green!10] (evo) at (-2.15,-5.7) {Time evolution\\$f(H)$};
	\node[qafnarr, fill=orange!10] (adi) at (2.15,-5.7) {Adiabatic path\\$H(s)$};
	\node[qafnarr, fill=yellow!18] (var) at (5.15,-5.7) {Variational\\hybrid};

	% ----- لایه ابزار -----
	\node[qafbox, text width=1.45cm, fill=green!6] (trot) at (-3.35,-7.35) {Trotter};
	\node[qafbox, text width=1.45cm, fill=green!6] (lcu) at (-2.15,-7.35) {LCU};
	\node[qafbox, text width=1.45cm, fill=green!6] (qsp) at (-0.95,-7.35) {QSP / QSVT};
	\node[qafbox, text width=1.7cm, fill=orange!12] (qpe) at (1.15,-7.35) {QPE};
	\node[qafbox, text width=1.55cm, fill=yellow!12] (vqe) at (5.15,-7.35) {VQE};

	% ----- لایه اجرا و خروجی -----
	\node[qafwide, fill=gray!14] (exec) at (0,-9.15) {Execution:\\circuit or analog schedule};
	\node[qafwide, fill=gray!8] (meas) at (0,-10.4) {Measurement / observable};
	\node[qafwide, fill=gray!8] (val) at (0,-11.55) {Classical validation};
	\node[qafwide, fill=gray!8] (res) at (0,-12.7) {Resource estimation};

	% برچسب لایه
	\node[lab] at (-6.55,-4.05) {QUESTION};
	\node[lab] at (-6.55,-5.7) {APPROACH};
	\node[lab] at (-6.55,-7.35) {TOOL};
	\node[lab] at (-6.55,-9.15) {RUN};

	% پس‌زمینه دو ستون پرسش
	\begin{scope}[on background layer]
		\node[qafpanel, fill=green!5, draw=green!40!black, fit=(dyn)(evo)(trot)(lcu)(qsp),
			inner xsep=5pt, inner ysep=6pt] {};
		\node[qafpanel, fill=orange!5, draw=orange!45!black, fit=(stat)(adi)(var)(vqe),
			inner xsep=5pt, inner ysep=6pt] {};
	\end{scope}

	% ----- اتصالات مسئله -----
	\draw[qafarr] (phys) -- (model);
	\draw[qafarr] (model) -- (H);
	\coordinate (hsplit) at (0,-3.15);
	\draw[qafarr] (H.south) -- (hsplit) -| (dyn.north);
	\draw[qafarr] (hsplit) -| (stat.north);

	% پرسش → رهیافت
	\draw[qafarr] (dyn) -- (evo);
	\draw[qafarr] (dyn.west) -- ++(-0.35,0) |- (ana.north);
	\draw[qafarr] (stat) -- (adi);
	\draw[qafarr] (stat.east) -- ++(0.35,0) |- (var.north);
	% پرسش طیفی می‌تواند از رهیافت تحول هم استفاده کند (QPE یا $f(H)$)
	\draw[qafarr, dashed] (stat.west) -- (evo.east)
		node[note, midway, above] {via $U(t)$};

	% رهیافت تحول → سه ابزار دیجیتال
	\draw[qafarr] (evo.south) -- ++(0,-0.18) -| (trot.north);
	\draw[qafarr] (evo) -- (lcu);
	\draw[qafarr] (evo.south) -- ++(0,-0.18) -| (qsp.north);

	% وردشی → VQE
	\draw[qafarr] (var) -- (vqe);

	% QPE از هر پیاده‌سازی U تغذیه می‌شود، نه فقط از یک ابزار
	\node[draw=green!45!black, densely dashed, rounded corners=2pt, inner sep=3pt,
		fit=(trot)(lcu)(qsp)] (utools) {};
	\draw[qafarr] (utools.east) -- (qpe.west)
		node[note, above, pos=0.55] {controlled $U$};

	% شینهٔ اجرا: همهٔ مسیرها جداگانه به اجرا می‌رسند؛ آنالوگ از گیت‌ها عبور نمی‌کند
	\coordinate (busL) at (-5.15,-8.35);
	\coordinate (busR) at (5.15,-8.35);
	\draw[thick, draw=black!55] (busL) -- (busR);
	\foreach \n in {ana,trot,lcu,qsp,qpe,adi,vqe} {
		\draw[qafarr] (\n.south) -- (\n.south |- busL);
	}
	\draw[qafarr] (0,-8.35) -- (exec.north);

	\draw[qafarr] (exec) -- (meas);
	\draw[qafarr] (meas) -- (val);
	\draw[qafarr] (val) -- (res);
\end{tikzpicture}
%
				}
		\end{latin}
	\caption{نقشه پژوهشی در چهار لایه پس از $H$: پرسش، رهیافت، ابزار، اجرا. کادر خط‌چین دور \lr{Trotter}/\lr{LCU}/\lr{QSP} یعنی هر سه پیاده‌سازی $U$ هستند و \lr{QPE} از همان گروه تغذیه می‌شود. پیکان خط‌چین از پرسش طیفی به \lr{QPE} نشان می‌دهد این ابزار پرسش طیفی را جواب می‌دهد ولی از لایه رهیافت تحول عبور می‌کند. مسیر آدیاباتیک دو بستر دارد: زمان‌بندی آنالوگ مستقیم، یا \lr{digital adiabatic} (پیکان خط‌چین به تروتر). آنالوگِ مهندسی $H$ بدون گیت دیجیتال به اجرا می‌رسد.}
	\label{fig:research-map}
\end{figure}

اتصالات شکل~\ref{fig:research-map} باید با همین تفکیک خوانده شوند:
\begin{itemize}
	\item \lr{Trotter} / \lr{LCU} / \lr{QSP} ابزارِ رهیافت تحول / $f(H)$ هستند؛ $U(t)=e^{-iHt}$ فقط یک کاربرد است. هر سه، جدا از هم، به اجرا می‌رسند.
	\item \lr{QPE} ابزار طیفی است اما ورودی‌اش تحول کنترل‌شده است؛ بنابراین از جعبه ابزارهای $U$ تغذیه می‌شود، نه از مسیر آدیاباتیک و نه فقط از \lr{LCU}.
	\item مسیر آدیاباتیک $H(s)$ یک \emph{رهیافت} موازی با تحول و وردشی است، نه زیرروالِ ذاتیِ \lr{Trotter} و نه پیش‌نیاز \lr{QPE}. اجرای آن یا آنالوگ است (زمان‌بندی روی سخت‌افزار) یا دیجیتال (\lr{digital adiabatic}: تروتر روی $H(s_k)$؛ فصل~\ref{chap:evo-adi}).
	\item \lr{VQE} زیرِ رهیافت وردشی / ترکیبی است، نه تجزیه $e^{-iHt}$.
	\item مهندسی آنالوگ از پرسش دینامیکی می‌آید و بدون عبور از گیت‌های دیجیتال به اجرا می‌رسد (فصل~\ref{chap:hardware}).
\end{itemize}


\begin{table}[htbp]
	\centering
	\footnotesize
	\setstretch{1.15}
	\setlength{\tabcolsep}{4pt}
	\caption{از پرسش علمی به رهیافت و ابزار نوعی.}
	\label{tab:question-approach-tool}
	\begin{tabularx}{\textwidth}{@{}
				>{\raggedright\arraybackslash}p{4.8cm}
				>{\raggedright\arraybackslash}p{2.5cm}
				>{\raggedright\arraybackslash}p{3.9cm}
				>{\raggedright\arraybackslash}p{3.5cm}
		%		>{\raggedright\arraybackslash}p{1.1cm}
				@{}}
			\toprule
			\rowcolor{gray!14}
			\textbf{اگر می‌خواهید} & \textbf{پرسش} & \textbf{رهیافت} & \textbf{ابزار نوعی} 
	%		& \textbf{فصل} 
			\\
			\midrule
			$\lvert\psi(t)\rangle$ یا مشاهده‌پذیر وابسته به \lr{t}
			& دینامیک
			& تحول زمانی / $f(H)$
			& \lr{Trotter}، \lr{LCU}، \lr{QSP}
	%		& \ref{chap:dyn-spec}، \ref{chap:digital-tools}
			\\
			\rowcolor{gray!8}
			انرژی پایه یا گاف
			& ایستا / طیفی
			& وردشی، آدیاباتیک، یا \lr{QPE}
			& \lr{VQE}، مسیر $H(s)$، \lr{QPE}
	%		& \ref{chap:dyn-spec}، \ref{chap:evo-adi}، \ref{chap:spectral-tools}
			\\
			ویژه مقدارهای $H$ با دقت بالا
			& طیفی
			& تحول کنترل‌شده $+$ برآورد فاز
			& شبیه‌سازی $U$ سپس \lr{QPE}
	%		& \ref{chap:spectral-tools}
			\\
			\rowcolor{gray!8}
			تابع پاسخ / \lr{Green}
			& پل دینامیک--طیف
			& $f(H)$ یا تحول
			& \lr{QSP}، \lr{QPE}
	%		& \ref{chap:dyn-spec}
			\\
			مهندسی مستقیم $H$ روی سخت‌افزار
			& دینامیک یا ایستا
			& آنالوگ
			& نگاشت $H_{\mathrm{target}}$
	%		& \ref{chap:hardware}
			\\
			\rowcolor{gray!8}
			ماندن در حالت پایه لحظه‌ای
			& ایستا
			& آدیاباتیک
			& زمان‌بندی $H(s)$ (آنالوگ یا دیجیتال)
	%		& \ref{chap:evo-adi}
			\\
			تعادل / مشاهده‌پذیر گرمایی
			& گرمایی
			& آماده‌سازی حالت گرمایی یا نمونه‌برداری
			& آنزات گرمایی، بازپخت، یا $f(H)$
	%		& 
			\\
			\rowcolor{gray!8}
			دینامیک با اتلاف / ناهمدوسی
			& سامانه باز
			& معادله مادر / مسیرهای کوانتومی
			& رقیق‌سازی، نقشه کراوس، شبیه‌سازی کانال
	%		& 
			\\
			جریان، ترابرد، هدایت
			& ترابرد
			& تحول یا پاسخ خطی
			& quench جریان، \lr{QSP}/پاسخ، اندازه‌گیری جریان
	%		& \ref{chap:dyn-spec}
			\\
			\bottomrule
		\end{tabularx}
\end{table}

\section{گسترش نقشه: گرمایی، سامانه باز، ترابرد}
\label{sec:map-extensions}

سطرهای جدول~\ref{tab:question-approach-tool} فقط دو پرسشِ بسته/طیفی را پوشش نمی‌دهند. سه خانواده دیگر که در پژوهش شبیه‌سازی زیاد دیده می‌شوند تا احتمالا بعداً با گسترش یابند.

\paragraph{مسائل گرمایی.}
پرسش: میانگین یک مشاهده‌پذیر در تعادل دمایی $T$، یا نمونه‌برداری از $\rho\sim e^{-\beta H}/Z$. این نه دینامیکِ خالصِ یک حالت خالص است و نه لزوماً فقط انرژی پایه. رهیافت‌های رایج: آماده‌سازی تقریبی حالت گرمایی (وردشی گرمایی، بازپخت)، یا محاسبه توابعی از $H$ که به $\beta$ وابسته‌اند. ابزارها با فصل‌های تحول و وردشی هم‌پوشانی دارند، اما خروجی اغلب $\langle O\rangle_\beta$ است نه $\lvert\psi(t)\rangle$.

\paragraph{سامانه کوانتومی باز.}
وقتی سامانه با محیط جفت است، تحول یکانیِ $e^{-iHt}$ کافی نیست؛ توصیف نوعی معادله مادر\footnote{\lr{master equation}.} (مثلاً لینبلادی) یا مسیرهای کوانتومی\footnote{\lr{quantum trajectories}.} است. روی نقشه: پرسش هنوز «دینامیک» است، اما رهیافت از یکانیِ بسته به کانال کاملاً مثبت / اتساع با محیط کمکی جابه‌جا می‌شود. ابزار دیجیتال می‌تواند همان کانال را شبیه‌سازی کند؛ آنالوگ گاهی اتلاف را ذاتاً دارد --- که هم فرصت است و هم نویز.

\paragraph{ترابرد.}
پرسش نوعی: جریان بار/اسپین، هدایت، یا پخش پس از یک گرادیان یا \lr{quench}. از نظر نقشه اغلب زیرِ دینامیک یا مرز پاسخ خطی می‌نشیند (تابع گرین، هدایت در فرکانس صفر). رهیافت: تحول زمانیِ مشاهده‌پذیر جریان، یا $f(H)$ برای پاسخ. همان ابزارهای فصل~\ref{chap:digital-tools} و پرسش‌های فصل~\ref{chap:dyn-spec}، با مشاهده‌پذیر و پروتکل اندازه‌گیریِ متفاوت.

\paragraph{دو بستر آدیاباتیک در خودِ نمودار.}
رهیافت $H(s)$ یکی است؛ اجرا دو گونه است (شکل~\ref{fig:research-map}):
\begin{itemize}
	\item \textbf{آدیاباتیک آنالوگ:} پارامترهای سخت‌افزار را مطابق $s(t)$ جارو می‌کنید؛ زمان‌بندی خودِ الگوریتم است (بخش~\ref{sec:analog}).
	\item \textbf{آدیاباتیک دیجیتال (\lr{digital adiabatic}):} در نقاط گسسته $s_k$ هامیلتونی $H(s_k)$ را با تروتر (یا ابزار $U$ دیگر) می‌سازید؛ پیکان خط‌چین از مسیر آدیاباتیک به جعبه \lr{Trotter} در نقشه همین را می‌گوید. \lr{QAOA} خویشاوند وردشیِ همین گسسته‌سازی است (بخش~\ref{sec:qaoa}).
\end{itemize}
پس «آدیاباتیک» را با «آنالوگ» یکی نگیرید: اولی رهیافت است، دومی یکی از بسترهای اجرا.


\begin{remark}[جای توسعه]
فصل یا زیربخش مستقل برای معادله مادر و مسیرهای کوانتومی؛ آماده‌سازی حالت گرمایی و پیچیدگی نمونه‌برداری؛ مثال عددی ترابرد روی زنجیره اسپینی کوچک؛ رسم جداگانه فقط‌آدیاباتیک با دو ستون آنالوگ/دیجیتال در صورت شلوغ شدن نقشه اصلی.
\end{remark}